\section{Conclusion}
\label{sec:conclusion}

This paper presented a systematic literature review of 54 studies on collaborative SLM--LLM systems across hybrid, multi-agent, ensemble, and retrieval-centered architectures, with monolithic LLMs treated as baselines rather than as a primary study class. The four research questions yield consistent, actionable answers.

\textbf{RQ1 (Performance).} Hybrid and multi-agent architectures consistently outperform monolithic SLM baselines and can match or exceed LLM baselines in narrower, structured settings. Performance gains are driven by orchestration strategy---selective escalation, role decomposition, verification loops---rather than by model size alone, and are strongest in domain-specific tasks with well-defined subtask boundaries.

\textbf{RQ2 (Efficiency).} The strongest efficiency gains, often exceeding 80\% cost reduction, arise from token-level routing and selective escalation. Latency effects are polarized: local SLM inference is fast, but orchestration layers frequently introduce overhead that negates throughput gains. Energy is documented in only 1 of the 54 included studies (S4); that single-study evidence shows that generation length and reasoning depth matter as much as model size for per-query energy use, and pushing small models into long Chain-of-Thought traces can exceed the per-query energy cost of a larger model prompted more directly~\cite{energytoken2025}. The remaining efficiency-adjacent figures in the corpus are compute or token-count proxies, such as the reduced computational redundancy reported in S29~\cite{resiliomate2025}, and should not be read as verified energy savings.

\textbf{RQ3 (Orchestration).} Orchestration mechanisms cluster into five recurring families: token-level verification, query-level routing, RAG, task decomposition, and multimodal fusion. The choice among them is determined primarily by deployment constraints and data modalities rather than by raw model capacity.

\textbf{RQ4 (Domain Suitability).} Domain-specific constraints---privacy mandates in healthcare, latency budgets in telecom, hardware limits at the edge, and syntax rigidity in code tasks---are the primary drivers of architecture selection. Collaborative SLM--LLM systems deliver their strongest results when the workflow can be partitioned to align each component with the constraint it best addresses.

Across all four questions, a consistent principle emerges: the presence of a \emph{well-calibrated} explicit control layer---router, planner, retriever, or fusion module---is the most reliable differentiator between high- and low-performing collaborative designs (39 of 54 included studies). Because the eligibility criteria required every included study to already employ some control mechanism, this differentiator concerns how well the layer is calibrated to the task, not whether one exists (Section~\ref{sec:cross_rq}).

\textbf{Design guidelines.} The following guidelines represent the most reliable actionable conclusions of this review, each supported by convergent evidence from at least four independent primary studies reporting consistent positive outcomes under comparable operational constraints.

\begin{enumerate}
    \item \textbf{G1 --- Prioritise well-calibrated control layers.} When choosing between monolithic and collaborative designs, deploy a routing, retrieval, or planning component as the primary architectural decision, and invest in calibrating its thresholds to the task; 39 of 54 studies attribute performance gains to a well-calibrated control layer rather than to model scale alone.
    \item \textbf{G2 --- Match mechanism to deployment constraint.} Use token-level routing when API cost or latency is the binding constraint; use RAG when hallucination reduction or data privacy is primary; use role-specialised multi-agent decomposition when reasoning depth across subtasks is required. Figure~\ref{fig:decision-flowchart} provides a practitioner-facing decision procedure.
    \item \textbf{G3 --- Report latency and cost alongside accuracy.} Accuracy-only evaluation is insufficient to establish deployment suitability; efficiency trade-offs are equally decision-relevant and currently underreported across the corpus.
    \item \textbf{G4 --- Prefer narrow, structured domains for initial hybrid deployment.} Collaborative SLM--LLM systems show their most consistent gains in healthcare, telecom, edge/IoT, and code tasks with well-defined subtask boundaries; performance in broad, open-ended generation tasks is less reliable and should be validated against a monolithic baseline before adoption.
\end{enumerate}

The literature is stronger in architectural innovation than in evaluation discipline. Many studies report accuracy improvements, but far fewer report latency, cost, or energy in a form that supports fair comparison. The current evidence is sufficient to show that hybridization \textit{can} work, but often insufficient to show \textit{when} it is the better deployment choice under realistic resource constraints. The strongest positive results remain concentrated in narrower or more structured settings such as telecom question answering, selected healthcare tasks, optimisation, and industrial monitoring. Future work should treat energy as a first-class evaluation outcome and report energy-per-token, reasoning length, hardware configuration, and escalation behavior under comparable conditions across multiple studies, not just the single S4 data point available today. Standardized benchmarks for collaborative SLM--LLM systems would further reduce the current reliance on paper-specific evaluation protocols.

The full comparative analysis table and data extraction spreadsheet supporting this review are available at \url{https://bit.ly/SLM-LLM-Data-2026}.
